\section{Conclusão}
\label{sec:conclusao}

Foi implementado e avaliado um sistema distribuído com eleição Bully, detector de
falhas por heartbeat e comunicação gRPC entre processos independentes em contêineres
Docker. A configuração por Docker Compose torna os cenários reprodutíveis,
enquanto os deadlines explícitos e os identificadores de eleição tornam observável
o tratamento de falhas e de mensagens obsoletas.

Os testes e os experimentos confirmaram que o maior nó ativo é escolhido como
coordenador em 100\% dos cenários avaliados. A redução do intervalo de heartbeat
de 3 s para 1 s reduziu o tempo médio de convergência de 4,900 s para 1,186 s.
A extensão para clusters com $N \in \{3, 5, 7, 10\}$ nós confirmou a curvatura
quadrática $\mathcal{O}(N^2)$ de mensagens do algoritmo Bully e demonstrou a
estabilidade do cluster sob falhas simultâneas de múltiplos líderes e falhas
intermediárias.

Como trabalhos futuros, sugere-se a avaliação em redes de longa distância (WAN)
com enlaces assimétricos e a incorporação de mecanismos de quórum majoritário
para mitigar partições de rede (\textit{split-brain}).
